\section{Construction requirements and remaining questions}\label{sec:roadmap}
These bounds distinguish solving a row for its center from solving it for
any participating coordinate. A construction on the prescribed supports
that exceeds the matching bound must violate the latter property. The paid
versions quantify the required failure: inverse fibres, conditional
marginals, or singular-block kernels must contribute enough along the fixed
schedule. Invertible changes of alphabet at individual coordinates do not
remove an information-set bound.

For linear codes, exceeding half the total coordinate-image dimension
also requires some rows with at least two components of rank greater than
one. The span of those exceptional coordinates must pay the excess.
This criterion permits arbitrarily large original coordinate widths and
does not require compatible scalar bases at different ports.

Independently, the squared recovery operator must have enough
off-diagonal rank: a full kernel of dimension $(1/2+\varepsilon)N$
requires $\rank\operatorname{offdiag}(W^2)\geq\varepsilon N/6$.
The exact loss on the kernel can be smaller than this global rank
payment. The quadratic identity further excludes every block-diagonal
square, over all characteristics. Square-covered support graphs are
already subject to a weighted polymatroid half bound, so the remaining
possibility must use off-square edges with nonzero two-step transmission.

The main transposition question (T) concerns a fixed finite-tape machine
running in $O(L)$ time, where $L=d^2b(d)$ and
$b(d)=\max(1,\lceil\log_2d\rceil)$. Nothing here supplies the necessary
subcubic representation for arbitrary tape computations, proves local
reversibility there, or pays the cost of implementing the inverse maps.
Those are separate construction or impossibility obligations. Even on a
triangle-free subcubic graph, arbitrary nonlinear recovery and unrestricted
singular-block recovery remain outside the unqualified half bound.
The next useful construction would simultaneously evade the reversible,
one-higher-rank, and small-defect criteria while meeting native tape costs.
A lower-bound approach must instead show why an adequate restriction or
payment is forced by the actual computation.
